% ============================================================================
%  A/04-tu-nguyen-va-cau-tao — file chương
%  Ngày tạo: 2026-09-06 | Sửa gần nhất: 2026-09-06 | Ôn gần nhất: chưa
%  Dependency: A/02. Mức độ: cốt lõi ở phần hậu tố và đoán nghĩa; tham khảo ở bảng gốc từ.
% ============================================================================
\documentclass[../../english.tex]{subfiles}
\begin{document}
\chapter{Từ nguyên và cấu tạo từ (Word Parts and Where They Came From)}
\ptrtarget{A/04}\ptrtarget{A/04-tu-nguyen-va-cau-tao}
\dependency{\ptr{A/02} cho khái niệm họ từ và gánh nặng học. Chương này giảm gánh nặng đó bằng cách chỉ ra rằng phần lớn từ kỹ thuật được lắp từ các mảnh có nghĩa cố định.}

\begin{tomtat}
Từ vựng khoa học tiếng Anh phần lớn không phải học thuộc mà là \emph{đọc được}: nó được lắp từ các gốc Latin và Hy Lạp cùng một tập nhỏ tiền tố và hậu tố. Chương này dạy ba việc: đọc hậu tố để biết từ loại và do đó biết vị trí của nó trong câu; đọc tiền tố và gốc để đoán nghĩa; và nhận ra ba tầng từ vựng của tiếng Anh --- lớp Germanic đời thường, lớp Pháp trung tính, lớp Latin và Hy Lạp trang trọng --- vì chính ba tầng này quyết định văn vực. Chương cũng nói rõ khi nào từ nguyên đánh lừa và khi nào nên ngừng đoán mà tra. Nếu chỉ nhớ một điều: học một gốc từ tốt cho bạn hàng chục từ, còn học một từ chỉ cho bạn một từ.
\end{tomtat}

\begin{nentang}
\kn{hình vị (morpheme)}{đơn vị nhỏ nhất mang nghĩa. \emph{Unbelievable} có ba: \emph{un-} (phủ định), \emph{believe} (gốc), \emph{-able} (có thể).}
\kn{gốc từ (root)}{phần mang nghĩa lõi, thường có nguồn Latin hoặc Hy Lạp trong từ vựng khoa học: \emph{spect} (nhìn), \emph{struct} (xây), \emph{graph} (viết).}
\kn{tiền tố (prefix)}{đứng trước gốc, thường đổi nghĩa mà không đổi từ loại: \emph{re-}, \emph{un-}, \emph{sub-}, \emph{over-}.}
\kn{hậu tố (suffix)}{đứng sau gốc, thường \emph{đổi từ loại}: \emph{-tion} tạo danh từ, \emph{-ise} tạo động từ, \emph{-ous} tạo tính từ.}
\kn{biến hình so với phái sinh}{biến hình chỉ đổi dạng ngữ pháp (\emph{show, shows, showed}); phái sinh tạo từ mới, có thể khác từ loại (\emph{show} \(\rightarrow\) \emph{showy}).}
\kn{bạn giả (false friend)}{từ trông giống một từ quen trong tiếng khác nhưng nghĩa khác: \emph{actually} không phải ``hiện tại'', \emph{eventually} không phải ``có thể''.}
\end{nentang}

\begin{khoigoi}
\h{Vì sao đọc được hậu tố lại giúp bạn hiểu câu, chứ không chỉ giúp đoán nghĩa từ?}
\h{Tiếng Anh có ba từ cho cùng một khái niệm: \emph{kingly}, \emph{royal}, \emph{regal}. Vì sao lại thừa như vậy, và sự thừa đó dùng để làm gì?}
\h{Khi gặp một từ kỹ thuật chưa từng thấy, bạn nên làm gì trước: tra hay đoán?}
\h{Từ nguyên đánh lừa ở những chỗ nào, và làm sao biết mình đang bị lừa?}
\h{Vì sao học một gốc từ lại rẻ hơn nhiều so với học từng từ riêng lẻ?}
\end{khoigoi}

Hãy thử một thí nghiệm nhỏ. Đọc từ \emph{photogrammetric} --- có thể bạn chưa từng gặp. Nhưng nếu biết \emph{photo-} là ánh sáng, \emph{-gram-} là cái được ghi lại, \emph{-metr-} là đo, và \emph{-ic} là hậu tố tính từ, thì bạn đoán được: thuộc về việc đo đạc dựa trên hình ảnh ghi lại được. Với một người làm thị giác máy tính, phỏng đoán đó đủ chính xác để đọc tiếp mà không đứt mạch.

Đây không phải trò chơi từ nguyên. Nó là một kỹ năng có giá trị thực dụng lớn vì hai lý do. Thứ nhất, từ vựng khoa học tiếng Anh được \emph{thiết kế} theo lối lắp ghép: khi cần đặt tên một khái niệm mới, người ta ghép các mảnh Latin hoặc Hy Lạp có nghĩa sẵn, nên nghĩa của từ mới suy được từ các mảnh. Thứ hai, việc đoán được nghĩa cho phép bạn đọc liên tục thay vì dừng tra --- và như \ptr{A/02} đã nói, giữ được mạch đọc là điều kiện để gặp lại từ đủ nhiều lần.

\section{Hậu tố: đọc từ loại trước khi đọc nghĩa}

Đây là phần có giá trị tức thì nhất, và nó ít được dạy hơn nhiều so với giá trị của nó. Hậu tố trong tiếng Anh gần như luôn báo \emph{từ loại}, mà biết từ loại là biết vị trí của từ trong câu --- tức là bạn phân tích được cấu trúc câu (\ptr{A/03}) ngay cả khi chưa biết nghĩa của từ.

\begin{center}\footnotesize
\begin{tabularx}{\linewidth}{@{}L{2.4cm}L{3.0cm}L{4.2cm}Y@{}}
\toprule
\textbf{Từ loại} & \textbf{Hậu tố hay gặp} & \textbf{Ví dụ} & \textbf{Ghi chú khi đọc}\\
\midrule
Danh từ & \emph{-tion, -sion, -ment, -ance, -ence, -ity, -ness, -ism} & \emph{estimation, alignment, robustness, accuracy} & Thường là danh hoá: hãy hỏi động từ gốc và ai thực hiện\\
Danh từ chỉ người/vật & \emph{-er, -or, -ist, -ant} & \emph{estimator, detector, analyst} & Chỉ tác nhân thực hiện hành động của gốc\\
Động từ & \emph{-ise/-ize, -ify, -ate, -en} & \emph{optimise, quantify, calibrate, strengthen} & Gần như luôn là ``làm cho trở thành''\\
Tính từ & \emph{-ous, -al, -ive, -able/-ible, -ic, -less} & \emph{robust\emph{ness} vs \emph{continuous}, iterative, tractable} & \emph{-able} nghĩa ``có thể được''; \emph{-less} là phủ định\\
Trạng từ & \emph{-ly} & \emph{significantly, roughly} & Bổ nghĩa cho động từ, tính từ, hoặc cả câu\\
\bottomrule
\end{tabularx}
\end{center}

Hai mẹo dùng ngay. \emph{Mẹo một:} khi gặp từ lạ, xác định hậu tố trước --- bạn lập tức biết nó là danh từ, động từ hay tính từ, và câu tự sáng ra một nửa. \emph{Mẹo hai:} khi viết, nếu câu của bạn có quá nhiều từ kết thúc bằng \emph{-tion} và \emph{-ment}, đó là dấu hiệu danh hoá quá nặng và cần viết lại theo \ptr{C/11}.

\section{Ba tầng từ vựng của tiếng Anh}

Đây là câu trả lời cho câu hỏi mở đầu thứ hai, và nó giải thích một đặc điểm của tiếng Anh khiến người học bối rối: vì sao có quá nhiều từ đồng nghĩa.

Lý do nằm ở lịch sử. Tiếng Anh có lõi Germanic --- các từ ngắn, đời thường, dùng hằng ngày. Sau cuộc chinh phục của người Norman năm 1066, tầng lớp cai trị nói tiếng Pháp trong nhiều thế kỷ, và hàng nghìn từ gốc Pháp đi vào tiếng Anh ở các lĩnh vực luật pháp, chính quyền, ẩm thực, nghệ thuật. Rồi từ thời Phục hưng, giới học giả mượn ồ ạt từ Latin và Hy Lạp để đặt tên cho các khái niệm khoa học và trừu tượng.

Kết quả là tiếng Anh thường có \emph{hai hoặc ba} từ cho cùng một khái niệm, mỗi từ ở một tầng, và chúng khác nhau về \emph{văn vực} chứ không về nghĩa quy chiếu.

\begin{tuvung}[Ba tầng của cùng một ý, khác nhau về văn vực]
\tv{ask}{hỏi}{Germanic; đời thường, trung tính, dùng khắp nơi}{I asked the authors for their code.}
\tv{question}{chất vấn}{Pháp; trung tính, có sắc thái nghi ngờ}{We question the assumption that noise is Gaussian.}
\tv{interrogate}{tra hỏi kỹ}{Latin; rất trang trọng hoặc kỹ thuật}{The tool interrogates the log for anomalies.}
\tv{start}{bắt đầu}{Germanic; đời thường}{We start with a simple baseline.}
\tv{begin}{bắt đầu}{Germanic, hơi trang trọng hơn \emph{start}}{The section begins with a definition.}
\tv{commence}{khởi sự}{Pháp/Latin; trang trọng, dễ thành khoa trương}{The trial commenced in March. (hiếm khi cần trong bài báo)}
\tv{buy}{mua}{Germanic; đời thường}{We bought two sensors.}
\tv{purchase}{mua}{Pháp; trang trọng, thiên hành chính}{Equipment was purchased under grant X.}
\end{tuvung}

Quy tắc thực dụng rút ra: \emph{từ gốc Germanic ngắn thường là lựa chọn tốt hơn trong văn khoa học hiện đại}, trừ khi từ Latin mang nghĩa kỹ thuật chính xác hơn. Người học hay nghĩ ngược lại --- rằng từ dài, gốc Latin mới ``học thuật'' --- và đó là nguyên nhân của thứ văn phong nặng nề mà mọi sách viết tốt đều cảnh báo (\ptr{C/11}).

\begin{trucgiac}
Người Việt có sẵn một mô hình rất sát để hiểu hiện tượng này: quan hệ giữa từ thuần Việt và từ Hán--Việt. ``Đàn bà'' và ``phụ nữ'' cùng chỉ một đối tượng nhưng khác văn vực; ``chết'', ``mất'', ``từ trần'' xếp thành một thang trang trọng. Tiếng Anh có đúng cấu trúc đó với ba tầng Germanic--Pháp--Latin, và nó ra đời cũng vì cùng một lý do lịch sử: một lớp từ vựng ngoại lai đi vào qua tầng lớp có học và ở lại trong các văn cảnh trang trọng. Nếu bạn đã có cảm giác về khi nào dùng ``phụ nữ'' và khi nào dùng ``đàn bà'', bạn đã có sẵn trực giác cần thiết --- chỉ cần chuyển nó sang tiếng Anh.
\end{trucgiac}

\section{Gốc từ: học một lần, dùng hàng chục lần}

Danh sách dưới đây không nhằm để học thuộc. Cách dùng đúng là: khi gặp một từ lạ trong tài liệu của bạn, tách nó ra và tra xem gốc đó còn xuất hiện ở những từ nào khác mà bạn đã biết --- việc nối được từ mới vào các từ cũ chính là thứ làm giảm gánh nặng học.

\begin{center}\footnotesize
\begin{tabularx}{\linewidth}{@{}L{2.2cm}L{2.6cm}Y@{}}
\toprule
\textbf{Gốc} & \textbf{Nghĩa} & \textbf{Từ trong tài liệu kỹ thuật}\\
\midrule
\emph{struct} & xây, dựng & structure, construct, infrastructure, destructive\\
\emph{spect / spic} & nhìn & inspect, perspective, conspicuous, retrospective\\
\emph{duc / duct} & dẫn & conduct, induction, deduce, reduce, conductive\\
\emph{ject} & ném & project, trajectory, injection, rejection\\
\emph{fer} & mang & transfer, inference, differential, reference\\
\emph{pos / pon} & đặt & compose, deposit, component, superposition\\
\emph{grad / gress} & bước & gradient, progressive, degradation, regression\\
\emph{form} & hình & transform, formulation, uniform, deformation\\
\emph{graph / gram} & viết, ghi & graph, diagram, photogrammetry, cartography\\
\emph{metr / meter} & đo & metric, symmetry, parameter, odometry\\
\emph{scop} & quan sát & microscope, telescopic, spectroscopy\\
\emph{therm} & nhiệt & thermal, isotherm, thermodynamics\\
\emph{morph} & hình dạng & morphology, isomorphic, polymorphism\\
\emph{syn / sym} & cùng & synchronise, symmetric, synthesis, symbiosis\\
\emph{topo} & nơi chốn & topology, topographic, topological\\
\bottomrule
\end{tabularx}
\end{center}

Tiền tố thì ít hơn và đáng thuộc hơn, vì chúng lặp lại rất dày: \emph{re-} (lại), \emph{un-} và \emph{in-/im-} (phủ định), \emph{dis-} (tách, phủ định), \emph{sub-} (dưới), \emph{super-} và \emph{over-} (trên, quá), \emph{inter-} (giữa), \emph{intra-} (bên trong), \emph{trans-} (xuyên qua), \emph{co-} (cùng), \emph{pre-} và \emph{post-} (trước, sau), \emph{anti-} (chống), \emph{semi-} (nửa), \emph{multi-} và \emph{poly-} (nhiều), \emph{mono-} và \emph{uni-} (một), \emph{iso-} (bằng nhau), \emph{non-} (không).

Chú ý một điểm dễ nhầm: \emph{in-} có hai nghĩa hoàn toàn khác nhau --- phủ định trong \emph{incorrect}, nhưng ``vào trong'' trong \emph{input}, \emph{insert}. Ngữ cảnh phân biệt, và đây là ví dụ nhỏ cho luận điểm ở \ptr{A/01} rằng ngữ cảnh mới quyết định nghĩa.

\section{Quy trình đoán nghĩa từ lạ}

Câu hỏi mở đầu thứ ba hỏi nên tra hay đoán. Câu trả lời là một quy trình bốn bước, và việc tra nằm ở bước cuối.

\emph{Bước 1 --- xác định từ loại qua hậu tố và vị trí trong câu.} Nhiều khi chỉ cần thế là đủ để đọc tiếp: nếu bạn biết đó là một tính từ bổ nghĩa cho \emph{estimator}, bạn không cần nghĩa chính xác để nắm cấu trúc câu.

\emph{Bước 2 --- đọc ngữ cảnh.} Câu trước và câu sau thường đã định nghĩa hoặc gợi ý. Trong văn khoa học, thuật ngữ mới hầu như luôn được giải thích ở lần xuất hiện đầu --- tìm cụm \emph{that is}, \emph{i.e.}, \emph{defined as}, hoặc phần trong ngoặc.

\emph{Bước 3 --- tách hình vị.} Tiền tố, gốc, hậu tố. Ghép nghĩa lại và kiểm xem có khớp với ngữ cảnh không.

\emph{Bước 4 --- tra, nhưng chỉ khi cần.} Tra khi từ đó xuất hiện lặp lại, khi nó nằm ở chỗ then chốt của lập luận, hoặc khi phỏng đoán của bạn không khớp ngữ cảnh. Nếu từ chỉ xuất hiện một lần ở chỗ không quan trọng, đọc tiếp là lựa chọn đúng.

\begin{cambay}
Từ nguyên đánh lừa ở ba chỗ. \emph{Trôi nghĩa}: nghĩa hiện tại có thể đã rời xa gốc --- \emph{nice} từng nghĩa là ``ngu dốt'', \emph{decimate} gốc là ``giết một phần mười'' nhưng nay dùng cho ``tàn phá nặng''. \emph{Bạn giả}: \emph{actually} nghĩa là ``thực ra'' chứ không phải ``hiện tại''; \emph{eventually} là ``rốt cuộc'' chứ không phải ``có thể''; \emph{sensible} là ``hợp lý'' còn \emph{sensitive} mới là ``nhạy''. \emph{Gốc giống nhau nhưng không liên quan}: \emph{universal} và \emph{university} có chung gốc nhưng \emph{universe} không giúp bạn đoán nghĩa của \emph{university} trong dùng hiện đại. Quy tắc an toàn: dùng từ nguyên để \emph{đoán}, dùng ngữ cảnh để \emph{kiểm}, và tin ngữ cảnh khi hai bên mâu thuẫn.
\end{cambay}

\section{Gốc rễ: vì sao tiếng Anh có ba tầng từ vựng}

Câu chuyện lịch sử ở mục 2 đáng kể chi tiết hơn, vì nó giải thích không chỉ hiện tượng đồng nghĩa mà cả cảm giác ``văn học thuật tiếng Anh nghe nặng''.

Tiếng Anh cổ là một ngôn ngữ Germanic, họ hàng gần với tiếng Đức và tiếng Hà Lan. Từ vựng lõi của nó vẫn còn nguyên tới hôm nay: các từ chỉ thân thể, gia đình, thời tiết, hành động hằng ngày --- \emph{hand, mother, rain, eat, go} --- gần như đều gốc Germanic, ngắn và không cần học.

Năm 1066, người Norman chiếm nước Anh và tiếng Pháp trở thành ngôn ngữ của triều đình, toà án và tầng lớp trên trong khoảng ba thế kỷ. Kết quả là một sự phân tầng để lại dấu vết tới hôm nay: con vật ngoài đồng mang tên Germanic (\emph{cow}, \emph{pig}, \emph{sheep}) còn món thịt trên bàn ăn của giới quý tộc mang tên Pháp (\emph{beef}, \emph{pork}, \emph{mutton}). Toàn bộ từ vựng pháp lý và hành chính tiếng Anh cho tới nay vẫn nặng gốc Pháp.

Tầng thứ ba đến từ thời Phục hưng, khi giới học giả mượn trực tiếp từ Latin và Hy Lạp để đặt tên các khái niệm mới. Có cả một giai đoạn tranh cãi trong thế kỷ 16 về việc vay mượn quá đà --- những từ mượn cầu kỳ bị chế giễu, và nhiều từ trong số đó đã biến mất, nhưng phần lớn thì ở lại. Về sau, khi khoa học hiện đại cần hàng vạn thuật ngữ mới, người ta tiếp tục dùng Latin và Hy Lạp làm nguồn --- không phải vì hoài cổ mà vì hai lý do thực dụng: các gốc đó có nghĩa ổn định và không xung đột với từ đời thường, và chúng ghép được với nhau theo quy tắc.

Ba hệ quả cho người học hôm nay. Thứ nhất, từ vựng khoa học có tính hệ thống cao và đoán được --- đó là toàn bộ nội dung mục 3 và 4. Thứ hai, tiếng Anh có nhiều cặp đồng nghĩa khác văn vực, và điều khiển văn vực là một kỹ năng viết thật sự (\ptr{D/16}). Thứ ba, và ít người nói ra: người Việt học từ khoa học tiếng Anh có lợi thế mà người bản ngữ không có --- bạn đã quen với việc một khái niệm có tên Hán--Việt trang trọng bên cạnh tên thuần Việt, nên cơ chế phân tầng này không xa lạ.

\begin{gocre}
\gr{Trước 1066 --- tiếng Anh cổ}{Một ngôn ngữ Germanic với từ vựng đời thường.}{Lõi từ vựng ngắn, cụ thể, vẫn còn nguyên tới nay: \emph{hand, eat, rain, go}.}
\gr{1066 và ba thế kỷ sau --- người Norman}{Tiếng Pháp thành ngôn ngữ của triều đình và toà án.}{Hàng nghìn từ Pháp đi vào ở tầng trang trọng; để lại cặp \emph{cow/beef}, \emph{ask/question} và toàn bộ từ vựng pháp lý.}
\gr{Thế kỷ 15--17 --- Phục hưng}{Cần từ cho các khái niệm học thuật và khoa học mới.}{Vay mượn ồ ạt từ Latin và Hy Lạp; tầng thứ ba hình thành, tạo ra bộ ba \emph{kingly/royal/regal}.}
\gr{Thế kỷ 18--20 --- khoa học hiện đại}{Cần đặt tên cho hàng vạn khái niệm mới, nhất quán giữa các thứ tiếng.}{Dùng gốc Latin và Hy Lạp làm nguyên liệu ghép: nghĩa ổn định, không xung đột với từ đời thường, ghép được theo quy tắc.}
\gr{Thế kỷ 20 --- phong trào viết rõ ràng}{Văn hành chính và khoa học ngày càng nặng nề vì lạm dụng tầng Latin.}{Khuyến nghị ưu tiên từ ngắn gốc Germanic khi có thể --- nguyên tắc được nói kỹ ở \ptr{C/11}.}
\end{gocre}

\section{Liên hệ ngang}

\begin{lienhe}
\lh{Tiếng Việt}{Từ Hán--Việt so với thuần Việt}{Song song gần như hoàn hảo với Latin so với Germanic: cùng chỉ một khái niệm, khác văn vực, và lớp vay mượn đi vào qua tầng lớp có học. Nếu bạn đã có trực giác về ``phụ nữ'' so với ``đàn bà'', bạn có sẵn công cụ để cảm văn vực tiếng Anh.}
\lh{Hoá học và sinh học}{Danh pháp có quy tắc ghép}{Tên hợp chất hoá học và tên loài đều được lắp từ các mảnh có nghĩa cố định --- cùng nguyên lý với từ khoa học tiếng Anh. Học quy tắc ghép một lần thì đọc được hàng nghìn tên.}
\lh{Lập trình}{Quy ước đặt tên và tiền tố tên hàm}{Lập trình viên cũng ghép từ có nghĩa để đặt tên (\emph{isValid}, \emph{preprocess}, \emph{deserialize}). Đọc tên hàm dài chính là đọc hình vị --- và cũng đọc ngược từ phải sang trái như danh từ ghép ở \ptr{A/03}.}
\lh{Toán học}{Thuật ngữ ghép từ gốc Hy Lạp}{\emph{Isomorphism}, \emph{homeomorphic}, \emph{polynomial} --- đọc được gốc thì đoán được nghĩa và thấy được quan hệ họ hàng giữa các khái niệm.}
\lh{Học máy}{Tách từ thành đơn vị nhỏ hơn để xử lý từ hiếm}{Các mô hình ngôn ngữ hiện đại tách từ thành mảnh nhỏ vì cùng lý do với người học: phần đuôi phân bố có quá nhiều từ hiếm, và các mảnh thì tái sử dụng được (\ptr{A/02}).}
\lh{Lịch sử và khảo cổ}{Từ mượn như dấu vết tiếp xúc văn hoá}{Bản đồ các lớp từ mượn trong một ngôn ngữ kể lại lịch sử tiếp xúc của cộng đồng nói ngôn ngữ đó --- một cách đọc lịch sử qua từ vựng.}
\end{lienhe}

\begin{traloi}
\tl{Vì hậu tố báo từ loại, mà từ loại quyết định vai trò của từ trong câu. Nếu bạn biết \emph{-ation} là danh từ và \emph{-ise} là động từ, bạn xác định được đâu là chủ ngữ, đâu là động từ chính ngay cả khi chưa biết nghĩa --- tức là bạn phân tích được cấu trúc câu theo \ptr{A/03}. Rất nhiều lần, nắm cấu trúc là đủ để đọc tiếp, và nghĩa chính xác sẽ tự lộ ra ở câu sau.}
\tl{Vì lịch sử: lõi Germanic, cộng lớp Pháp sau năm 1066, cộng lớp Latin và Hy Lạp từ thời Phục hưng. Sự ``thừa'' đó không vô ích --- nó cho tiếng Anh một thang văn vực rất tinh: \emph{kingly} thân mật, \emph{royal} trung tính, \emph{regal} trang trọng. Người viết dùng thang này để điều chỉnh giọng, và người đọc nên đọc ra nó như một tín hiệu về mức trang trọng và về thái độ.}
\tl{Đoán trước, tra sau, theo bốn bước: xác định từ loại qua hậu tố; đọc ngữ cảnh xem có định nghĩa ngay không; tách hình vị và ghép nghĩa; rồi mới tra --- và chỉ tra khi từ lặp lại, nằm ở chỗ then chốt, hoặc phỏng đoán không khớp ngữ cảnh. Lý do ưu tiên đoán không phải để tiết kiệm thời gian tra mà để \emph{giữ mạch đọc}, vì đứt mạch làm bạn mất cả nội dung lẫn cơ hội gặp lại từ trong ngữ cảnh.}
\tl{Ba chỗ: nghĩa trôi theo thời gian (\emph{nice}, \emph{decimate}); bạn giả với tiếng khác hoặc với từ trông giống (\emph{actually}, \emph{eventually}, \emph{sensible} so với \emph{sensitive}); và các gốc trùng hình thức nhưng khác nguồn. Cách biết mình bị lừa rất đơn giản: phỏng đoán từ từ nguyên \emph{không khớp} với ngữ cảnh. Khi hai bên mâu thuẫn, luôn tin ngữ cảnh --- từ nguyên chỉ là công cụ tạo giả thuyết.}
\tl{Vì gánh nặng học phụ thuộc vào mức độ từ mới nối được vào những gì bạn đã biết. Một gốc như \emph{struct} xuất hiện trong hàng chục từ; khi biết gốc, mỗi từ mới trong họ chỉ còn là một tổ hợp gốc--phụ tố cần xác nhận, thay vì một chuỗi ký tự vô nghĩa cần nhớ. Đây cũng là lý do các con số ở \ptr{A/02} tính theo \emph{họ từ}: một khi biết gốc và các phụ tố phổ biến, cả họ trở nên gần như miễn phí.}
\end{traloi}

\begin{dongian}
Rất nhiều từ khoa học tiếng Anh không phải để học thuộc, mà để \emph{đọc ra}. Chúng được lắp lại từ những mảnh nhỏ có nghĩa cố định, giống như chữ Hán ghép từ các bộ.

Ví dụ \emph{odometry}: \emph{odo-} là đường đi, \emph{-metry} là phép đo --- phép đo quãng đường. \emph{Photogrammetry}: ánh sáng, cái được ghi lại, phép đo --- đo đạc dựa trên ảnh. Biết vài chục mảnh như vậy thì hàng trăm thuật ngữ trở nên đoán được, và bạn không phải dừng lại tra giữa chừng.

Phần cuối của từ còn cho một thông tin nữa, và nó hữu ích ngay cả khi bạn chưa đoán được nghĩa: nó cho biết từ đó là danh từ, động từ hay tính từ. \emph{-tion} là danh từ, \emph{-ise} là động từ, \emph{-ous} là tính từ. Chỉ cần biết vậy là bạn đã đọc được cấu trúc câu, và thường thì câu sau sẽ tự giải thích nghĩa.

Còn một điều đáng biết về tiếng Anh nói chung. Nó có ba lớp từ chồng lên nhau vì ba đợt lịch sử: lớp gốc bản địa (từ ngắn, đời thường), lớp tiếng Pháp (sau khi người Norman sang năm 1066), và lớp Latin--Hy Lạp (khi giới học giả đặt tên cho khoa học). Vì thế mới có ba từ cho cùng một ý, chỉ khác mức trang trọng --- y hệt tiếng Việt có ``đàn bà'' và ``phụ nữ''.

Điều này dẫn tới một lời khuyên trái với trực giác của nhiều người học: trong văn khoa học hiện đại, từ ngắn gốc bản địa thường \emph{tốt hơn} từ dài gốc Latin. Viết \emph{we use} thay vì \emph{we utilise} không làm bài của bạn kém học thuật hơn; nó chỉ làm câu nhẹ đi.
\motcau{Học một gốc cho bạn hàng chục từ; học một từ cho bạn đúng một từ.}
\chuadongian{Có những từ mà nghĩa hiện nay đã rời hẳn gốc; với chúng, từ nguyên chỉ gây nhầm và cách duy nhất là học theo cách dùng thật.}
\end{dongian}

\begin{onbai}
\ar[3]{Hậu tố báo gì}{Từ loại: \emph{-tion/-ment/-ity} danh từ, \emph{-ise/-ify/-ate} động từ, \emph{-ous/-al/-ive/-able} tính từ, \emph{-ly} trạng từ. Biết từ loại là đọc được cấu trúc câu.}
\ar[3]{Ba tầng từ vựng tiếng Anh}{Germanic (đời thường, ngắn), Pháp (trung tính tới trang trọng), Latin--Hy Lạp (trang trọng, kỹ thuật). Khác văn vực chứ không khác nghĩa quy chiếu.}
\ar[3]{Quy trình đoán nghĩa bốn bước}{Từ loại qua hậu tố \(\rightarrow\) ngữ cảnh (tìm \emph{i.e.}, \emph{defined as}) \(\rightarrow\) tách hình vị \(\rightarrow\) chỉ tra khi từ lặp lại hoặc ở chỗ then chốt.}
\ar[3]{Vì sao ưu tiên đoán}{Để giữ mạch đọc; đứt mạch làm mất cả nội dung lẫn cơ hội gặp lại từ trong ngữ cảnh (\ptr{A/02}).}
\ar[3]{Ba chỗ từ nguyên đánh lừa}{Nghĩa trôi (\emph{nice}, \emph{decimate}); bạn giả (\emph{actually}, \emph{eventually}, \emph{sensible/sensitive}); gốc trùng hình thức khác nguồn. Mâu thuẫn thì tin ngữ cảnh.}
\ar[2]{Vì sao học gốc rẻ hơn học từ}{Gánh nặng học giảm khi từ mới nối được vào cái đã biết; một gốc mở khoá cả một họ từ.}
\ar[2]{Lời khuyên trái trực giác về văn vực}{Trong văn khoa học hiện đại, từ ngắn gốc Germanic thường tốt hơn từ dài gốc Latin: \emph{use} chứ không \emph{utilise} (\ptr{C/11}).}
\ar[2]{Tiền tố đáng thuộc}{\emph{re-, un-, in-/im-, dis-, sub-, super-, over-, inter-, intra-, trans-, co-, pre-, post-, anti-, semi-, multi-, mono-, iso-, non-}.}
\ar[2]{Bẫy \emph{in-}}{Hai nghĩa khác hẳn: phủ định (\emph{incorrect}) và ``vào trong'' (\emph{input}); ngữ cảnh phân biệt.}
\ar[2]{Song song với tiếng Việt}{Latin/Germanic tương ứng Hán--Việt/thuần Việt: cùng cơ chế phân tầng văn vực do lịch sử vay mượn.}
\ar[1]{Vì sao khoa học dùng gốc Latin--Hy Lạp}{Nghĩa ổn định, không xung đột với từ đời thường, và ghép được theo quy tắc --- nên thuật ngữ mới đoán được nghĩa.}
\end{onbai}

\begin{hoidap}
\qa{Tôi có nên học thuộc bảng gốc từ không?}{Không nên học thuộc theo bảng. Cách hiệu quả hơn: mỗi khi gặp một từ lạ trong tài liệu \emph{của bạn}, hãy tách nó ra và tự hỏi ``gốc này còn nằm trong từ nào tôi đã biết''. Việc nối được từ mới vào mạng từ cũ mới là thứ tạo trí nhớ; đọc một bảng ngoài ngữ cảnh thì tối đa cho bạn cảm giác quen. Bảng ở mục 3 nên dùng như tài liệu tra, không phải danh sách phải học.}
\qa{Từ điển từ nguyên có hữu ích cho việc học không?}{Có, nhưng ở vai trò phụ. Nó hữu ích khi bạn muốn \emph{nhớ} một từ khó: câu chuyện nguồn gốc tạo một móc treo trong trí nhớ. Nó không hữu ích để \emph{xác định nghĩa hiện tại}, vì nghĩa trôi theo thời gian. Cách dùng hợp lý: tra nghĩa hiện tại ở từ điển học tập, và chỉ tra từ nguyên\cite{etymonline} khi từ đó khó nhớ hoặc khi bạn muốn thấy quan hệ họ hàng giữa các thuật ngữ.}
\qa{Làm sao biết một từ thuộc tầng nào để chọn cho đúng văn vực?}{Ba dấu hiệu nhanh. Độ dài và hình thức: từ ngắn, một hai âm tiết, không có phụ tố Latin thường thuộc lớp Germanic. Hậu tố: \emph{-tion, -ity, -ous} là dấu hiệu gốc Latin. Và nhãn trong từ điển học tập --- \emph{formal}, \emph{technical}. Khi vẫn không chắc, hãy tra tần suất theo thể loại trong kho ngữ liệu\cite{coca}: nếu từ xuất hiện chủ yếu ở văn bản học thuật thì nó thuộc lớp trang trọng.}
\qa{Trong bài báo, dùng từ dài có làm bài chuyên nghiệp hơn không?}{Không, và thường ngược lại. Người phản biện có kinh nghiệm đọc rất nhiều bài, và văn phong nặng nề làm chậm họ chứ không gây ấn tượng. Điều tạo cảm giác chuyên nghiệp là dùng \emph{đúng} thuật ngữ kỹ thuật của ngành --- nơi từ Latin là bắt buộc vì nó có nghĩa chính xác --- còn ở phần văn xuôi nối kết thì dùng từ đơn giản. Nói cách khác: chính xác ở thuật ngữ, giản dị ở phần còn lại.}
\qa{Tại sao tôi hay nhầm \emph{sensible} với \emph{sensitive}, \emph{economic} với \emph{economical}?}{Vì hai từ chung gốc nhưng hậu tố khác nhau tạo nghĩa khác nhau, và tiếng Việt dịch cả hai gần như một. Đây là nhóm đáng học riêng: \emph{sensible} (hợp lý) so với \emph{sensitive} (nhạy); \emph{economic} (thuộc kinh tế) so với \emph{economical} (tiết kiệm); \emph{classic} (kinh điển) so với \emph{classical} (cổ điển); \emph{historic} (có tính lịch sử, quan trọng) so với \emph{historical} (thuộc về lịch sử). Cách xử lý: ghi thành cặp trong sổ, mỗi từ kèm một câu thật.}
\qa{Từ vựng kỹ thuật của ngành tôi có nhiều từ ghép rất dài. Có mẹo đọc không?}{Có, và nó đã nói ở \ptr{A/03}: đọc ngược từ phải sang trái. Danh từ cuối cùng là cái chính, mọi thứ trước nó bổ nghĩa. Với \emph{visual inertial odometry initialisation failure}, đọc là: \emph{failure} \(\rightarrow\) của \emph{initialisation} \(\rightarrow\) của \emph{odometry} \(\rightarrow\) loại \emph{visual inertial}. Kết hợp với việc tách hình vị trong từng từ, bạn xử lý được gần như mọi cụm thuật ngữ dài mà không cần tra.}
\end{hoidap}

\begin{baitap}
\bt[1]{Lấy 10 thuật ngữ trong ngành bạn và tách chúng thành hình vị}{Tài liệu ngành}{Ghi nghĩa từng mảnh; kiểm xem nghĩa ghép có khớp nghĩa thật không.}
\bt[1]{Với mỗi hậu tố ở bảng mục 1, tìm ba từ trong tài liệu của bạn}{Tài liệu ngành}{Mục tiêu là nhận diện tự động, không phải học thuộc bảng.}
\bt[2]{Tìm 5 cặp từ cùng nghĩa khác tầng và xếp theo văn vực}{\cite{coca}}{Ví dụ \emph{begin/commence}, \emph{use/utilise}; kiểm bằng tần suất theo thể loại.}
\bt[2]{Chọn 5 gốc từ ở bảng mục 3 và liệt kê mọi từ bạn đã biết chứa gốc đó}{Trí nhớ của bạn}{Bài tập này nối từ mới vào mạng cũ --- chính là cơ chế giảm gánh nặng học.}
\bt[2]{Lập danh sách 10 cặp bạn giả gây nhầm cho bạn}{Nhật ký lỗi của bạn}{Mỗi từ kèm một câu thật; ôn lại sau một tuần.}
\bt[3]{Đọc một trang tài liệu mới và áp quy trình bốn bước cho mọi từ lạ}{Tài liệu ngành}{Ghi lại tỉ lệ đoán đúng; con số này thường cao hơn bạn nghĩ và tăng nhanh theo thời gian.}
\bt[3]{Viết lại một đoạn của bạn, thay các từ Latin dài bằng từ ngắn khi có thể}{Bài viết của bạn}{Giữ nguyên thuật ngữ kỹ thuật; chỉ đổi phần văn xuôi nối kết. So hai bản.}
\end{baitap}

\section*{Kết chương và đọc tiếp}

Chương này khép lại phần nền tảng bằng một công cụ rất thực dụng: đọc được cấu tạo từ để đoán nghĩa và để nhận biết văn vực. Hai điều đáng mang theo: hậu tố cho bạn từ loại --- đủ để phân tích câu ngay cả khi chưa biết nghĩa; và ba tầng từ vựng của tiếng Anh là một thang văn vực mà bạn đã có sẵn trực giác nhờ tiếng Việt.

Phần B chuyển sang việc đọc thật. Với bộ công cụ của phần A --- ba tầng nghĩa, chín mặt của một từ, cách tách câu, và cách đọc cấu tạo từ --- bạn đã đủ để bắt đầu làm việc với văn bản dài. Chương \ptr{B/05} bắt đầu bằng câu hỏi mà người đọc tài liệu chuyên ngành hay bỏ qua: \emph{lần đọc này của tôi nhằm mục đích gì}, và vì sao trả lời câu đó trước lại tiết kiệm được nhiều giờ.
\end{document}
